% answers_short.tex —— 2013 年·广东卷·理科 · 简短答案速查
% 由 tools/make_short_answers.py 自 answers.tex 的 \gksolhead/\gkanswer 提取，
% 并与 2013/2013_数学真题_广东理_答案要点.md 逐题比对后生成，勿手改（改源 MD 后重跑）。

\begin{gknotice}
\textbf{使用说明}：本节只给最终答案便于快速核对，完整推导见解析版。答案为编写者独立推导，\textbf{非官方标准答案}。
\end{gknotice}

\gkpart{选择题}
\begin{center}\begin{tabular}{ccccccccc}
\toprule
题号 & 1 & 2 & 3 & 4 & 5 & 6 & 7 & 8 \\
\midrule
答案 & D & C & C & A & B & D & B & B \\
\bottomrule
\end{tabular}\end{center}

\gkpart{填空题}
\begin{itemize}[leftmargin=2.6em]
  \item[\textbf{9.}] $\{x\mid -2<x<1\}$。$(x+2)(x-1)<0$，故 $-2<x<1$。
  \item[\textbf{10.}] $-1$。$y'=k+\dfrac1x$，在 $x=1$ 处 $k+1=0$。
  \item[\textbf{11.}] $7$。循环体为 $s=s+(i-1)$，依次累加 $0,1,2,3$ 到初值 $1$ 上。
  \item[\textbf{12.}] $20$。$a_3+a_8=2a_1+9d=10$，而 $3a_5+a_7=4a_1+18d=2(2a_1+9d)=20$。
  \item[\textbf{13.}] $6$。$T=\{(0,1),(0,4),(1,3),(2,2),(3,1),(4,0)\}$，其中五个最大值点在 $x+y=4$ 上，另一点与它们各连一条，共 $1+5=6$ 条。
  \item[\textbf{14.}] $\rho(\cos\theta+\sin\theta)=2$。圆为 $x^2+y^2=2$，在 $(1,1)$ 处的切线为 $x+y=2$。
  \item[\textbf{15.}] $2\sqrt3$。设 $BC=b$，$AC=h$，由 $AD=AB=6$、$AE=4$ 及 $E$ 在切线上得 $h^2=2b^2$，结合 $b^2+h^2=36$ 得 $b=2\sqrt3$。
\end{itemize}

\gkpart{解答题}
\begin{itemize}[leftmargin=2.6em]
  \item[\textbf{16.}] (1) $1$。$f\left(-\dfrac{\pi}{6}\right)=\sqrt2\cos\left(-\dfrac{\pi}{4}\right)=1$；(2) $\dfrac{17}{25}$。由 $\cos\theta=\dfrac35$、$\theta$ 在第四象限得 $\sin\theta=-\dfrac45$，再用二倍角公式计算
  \item[\textbf{17.}] (1) $22$。六个数据之和为 $132$，除以 $6$ 得 $22$；(2) $4$ 名。样本中大于 $22$ 的有 $25,30$ 两个，占 $\dfrac13$，故估计 $12\times\dfrac13=4$ 名；(3) $\dfrac{16}{33}$。$\dfrac{\mathrm C_4^1\mathrm C_8^1}{\mathrm C_{12}^2}=\dfrac{32}{66}=\dfrac{16}{33}$
  \item[\textbf{18.}] (1) 见解析版；(2) $\dfrac{\sqrt{15}}5$。建立空间直角坐标系，平面 $BCDE$ 与平面 $A'CD$ 的法向量分别为 $(0,0,1)$ 与 $(\sqrt3,-\sqrt3,-3)$，夹角余弦为 $\dfrac{3}{\sqrt{15}}=\dfrac{\sqrt{15}}5$
  \item[\textbf{19.}] (1) $4$。令 $n=1$，由 $2S_1=2$ 得 $2=a_2-2$，故 $a_2=4$；(2) $a_n=n^2$。相减得 $a_{n+1}/(n+1)=a_n/n+1$，故 $a_n/n=n$；(3) 见解析版
  \item[\textbf{20.}] (1) $x^2=4y$。由 $\dfrac{c+2}{\sqrt2}=\dfrac{3\sqrt2}{2}$ 得 $c=1$；(2) $x_0x-2y-2x_0+4=0$。由两条切线方程消去切点参数得极线（弦 $AB$），再代入 $y_0=x_0-2$；(3) $\dfrac92$。$|AF|\cdot|BF|=(y_1+1)(y_2+1)=x_0^2+(y_0-1)^2=2\left(x_0-\dfrac32\right)^2+\dfrac92$
  \item[\textbf{21.}] (1) 单调递增区间为 $(-\infty,0)$ 与 $(\ln2,+\infty)$，单调递减区间为 $(0,\ln2)$。$f'(x)=x(\mathrm e^x-2)$；(2) $M=(k-1)\mathrm e^k-k^3$。$f'(x)=x(\mathrm e^x-2k)$，$f$ 在 $\ln(2k)$ 处取极小值，最大值在端点 $0$ 或 $k$ 处，比较 $f(0)=-1$ 与 $f(k)$ 即得
\end{itemize}
